\chapter{Experimental evaluation}
\label{chap:evaluation}

\section{Methodology}
\label{sec:methodology}

\paragraph{Data.} Two kinds of instances are used. The \emph{benchmarks} are 579 public job shop
and flexible job shop instances converted into the unified schema
(Table~\ref{tab:benchmarks}); they have no due dates, so their only objective is the makespan,
and for 152 JSPLIB instances the best known makespan is available. The \emph{generated
instances} are 200 FACTORY and 200 PA instances produced by the generator of
Section~\ref{sec:generator} with seeds 0--199; they have due dates, job weights and random
objective weights.

\begin{table}[htpb]
  \centering
  \small
  \begin{tabular}{@{}lrr@{}}
    \toprule
    \textbf{Dataset} & \textbf{Instances} & \textbf{With best known value} \\
    \midrule
    Brandimarte~\cite{brandimarte1993}          & 10  & -- \\
    Hurink (edata, rdata, sdata, vdata)         & 264 & -- \\
    Barnes                                      & 21  & -- \\
    Dauzère-Pérès and Paulli                    & 18  & -- \\
    Fattahi                                     & 20  & -- \\
    Kacem                                       & 4   & -- \\
    JSPLIB (incl.\ Taillard~\cite{taillard1993}) & 162 & 152 \\
    Demirkol                                    & 80  & -- \\
    \midrule
    Total                                       & 579 & 152 \\
    \bottomrule
  \end{tabular}
  \caption{Benchmark instances converted into the unified schema}
  \label{tab:benchmarks}
\end{table}

\paragraph{Metrics.} On the benchmarks, the quality of a plan is measured by the gap of its
makespan from the best known value. On the generated instances it is measured by the weighted
tardiness and by the cost of Section~\ref{sec:objective}, which combines makespan, weighted
tardiness and setup with the objective weights of the instance. Every plan is verified by the
feasibility checker before it is counted.
% TODO [Gantt 10, P1, 30 Nov-11 Dec 2026 - Pilot]: metrics of M2: cost and weighted tardiness
% against ATC/EDD, share of instances where M2 wins, response time; test set of the dataset.
% TODO [Gantt 11, P1, 7-23 Dec 2026 - M1 context]: metrics of M1: MAE, coverage of the prediction
% intervals.

\section{Results}
\label{sec:results}

\subsection{Due dates of the generator}
\label{sec:eval-generator}

The first version of the generator set the due date of a job as its release date plus its
critical path plus a random slack, without considering the load of the shared resources. With
this rule the due dates were too loose: even with EDD, 61\% of the FACTORY instances and 79\% of
the PA instances had no tardiness at all, so tardiness could not discriminate between good and
bad schedules. Adding the estimated makespan to the slack term
(Section~\ref{sec:generator}) reduced these shares to 12\% and 18\%. Several slack ranges were
compared; the chosen ones (0--0.6 for FACTORY, 0--0.4 for PA) keep most instances tardy without
making all jobs hopelessly late.

\subsection{Dispatching rules}
\label{sec:eval-rules}

All 579 benchmark instances receive a feasible plan from both rules, in 1 to 105~ms per
instance. On JSPLIB the average gap from the best known makespan is 36.3\% for EDD and 48.7\% for
ATC ($K=2$), and no makespan falls below the known lower bound. The gap is large, as expected:
without due dates EDD reduces to ``earliest completion first'' and ATC to weight over duration,
and neither is designed for the makespan.

On the generated instances (Table~\ref{tab:rules}) ATC clearly beats EDD on weighted tardiness:
by 20\% on FACTORY and by 40\% on PA with $K=1$. The parameter $K=1$ is better than or equal to
$K=2$ in almost all configurations.

\begin{table}[htpb]
  \centering
  \small
  \begin{tabular}{@{}lrrr@{}}
    \toprule
    \textbf{Profile} & \textbf{EDD} & \textbf{ATC, $K=1$} & \textbf{ATC, $K=2$} \\
    \midrule
    FACTORY & 925  & 739  & 747 \\
    PA      & 4115 & 2465 & 2758 \\
    \bottomrule
  \end{tabular}
  \caption{Average weighted tardiness of the dispatching rules on the generated instances}
  \label{tab:rules}
\end{table}

\subsection{CP-SAT teacher}
\label{sec:eval-teacher}

As a check of the model, the teacher was run on the 22 JSPLIB instances with a proven optimum
among ft06, ft10, ft20 and la01--la19, with a limit of 5~seconds each: it proved optimality and
found the optimal makespan in all cases.

Table~\ref{tab:teacher} reports the results on the generated instances, with a limit of 10
seconds and 8 threads per instance.

\begin{table}[htpb]
  \centering
  \small
  \begin{tabular}{@{}lrrrr@{}}
    \toprule
    \textbf{Profile} & \textbf{Proven optimal} & \textbf{Median gap} & \textbf{Cost vs best rule}
      & \textbf{Avg.\ time (s)} \\
    \midrule
    FACTORY & 95 / 200  & 4.5\% & $-27.9\%$ & 5.9 \\
    PA      & 141 / 200 & 0.0\% & $-32.8\%$ & 3.9 \\
    \bottomrule
  \end{tabular}
  \caption{CP-SAT teacher on the generated instances. The gap is between the solution and the
           lower bound proven by CP-SAT; the best rule is the best of EDD, ATC with $K=1$ and ATC
           with $K=2$ on each instance.}
  \label{tab:teacher}
\end{table}

The teacher improves on the best dispatching rule by about 30\%, so there is room for a learned
policy to beat ATC. Most FACTORY instances that reach the time limit have sequence-dependent
setups: for them the lower bound proven by CP-SAT is weak, so a high gap does not necessarily
mean a poor solution. With several threads CP-SAT is not deterministic, and two runs can return
different plans when optimality is not proven; this is why the solutions are saved and the
dataset is built from them.

\subsection{Training dataset}
\label{sec:eval-dataset}

Table~\ref{tab:dataset} shows the size of the first version of the dataset for M2, built from
the 400 teacher solutions without any filter on the gap. A step has on average 7.4 candidates in
FACTORY and 12.8 in PA; only 4.5\% and 2.7\% of the steps have a single candidate, i.e.\ no
real decision.

\begin{table}[htpb]
  \centering
  \small
  \begin{tabular}{@{}lrrr@{}}
    \toprule
    \textbf{Split} & \textbf{Instances} & \textbf{Steps (queries)} & \textbf{Candidates} \\
    \midrule
    Training   & 287 & 16\,198 & 178\,550 \\
    Validation & 64  & 3\,472  & 35\,458 \\
    Test       & 49  & 2\,705  & 24\,170 \\
    \bottomrule
  \end{tabular}
  \caption{Dataset v1 for the policy M2}
  \label{tab:dataset}
\end{table}

The plan obtained by replaying the labels in the dispatcher costs on average 3.2\% more than the
teacher on FACTORY and 5.0\% more on PA. The loss comes from the dispatcher, which appends
activities to the resources and proposes only active candidates, so some sequences of the
teacher cannot be reproduced exactly; in 17 FACTORY instances the replay is even better than a
non-optimal teacher solution. The replay is still 26.0\% (FACTORY) and 30.3\% (PA) cheaper than
the best rule. This is the upper limit of what a perfect imitator can achieve with this
dispatcher, and the target for the policy M2.

% TODO [Gantt 10 / Milestone M1, 11 Dec 2026 - Pilot outcome]: results of M2 against EDD, ATC and
% CP-SAT on the test set; go / no-go decision.
% TODO [Gantt 11, P1, 7-23 Dec 2026 - M1 context]: results of M1 on the event logs.
% TODO [Gantt 14, P1, 8-29 Jan 2027 - Generalization experiments]: M2 on instances larger than
% the training ones, on the other profile, on the benchmarks (Brandimarte, JSPLIB) and with
% calibrated profiles.
% TODO [Gantt 29, P1, 18 Jan-3 Feb 2027 - Experimental evaluation and shadow mode]: response time
% of the AI service, plans computed in shadow mode next to the real ones.

\section{Discussion}
\label{sec:discussion}
% TODO: where the model works, where it does not, and why.
